\begin{definition}
Let $T$ be a tree-sequence, let $C\subseteq[k]$, and let $\varphi$ and
$\pi$ be partial $k$-edge-colorings.  The coloring $\pi$ is
$(T,C,\varphi)$-stable if both of the following hold:
\begin{enumerate}
\item for every edge identity $f$ incident with $T$, if $\varphi(f)$ is a
color in $\overline\varphi(T)\cup C$, then $\pi(f)=\varphi(f)$;
\item for every $v\in V(T)$,
$\overline\pi(v)=\overline\varphi(v)$.
\end{enumerate}
This is the manuscript's stable-coloring relation; in its applications both
colorings are proper and leave the same distinguished edge uncolored.
\end{definition}
